\section{Selecting tail tests}
\label{sec:selections}

The paired argument uses both parts of every relative potential. Sometimes
one needs a single test, or wants to decide separately which exit components
contribute to each endpoint. We give these choices explicitly. They preserve
more information about the witnessing test, although their numerical midpoint
bounds are weaker than Theorem~\ref{thm:four}. The later path and segment
arguments do not depend on this section; the numerical estimates below
illustrate what each selection of a witness retains.

Throughout this section, $K$ is one of the four norming sets of
Section~\ref{sec:coefficient-tests}, with coordinate set $I$.
Let $D$ be a finite initial set and let $f\in K$ be supported in
$D\cap I$. An exit component is a subtree whose first vertex lies
outside $D$ and whose parent $r$ lies in $D$. Its inherited head value
is $c=P_f(r)$. Given an arbitrary test $g\in K$, define its relative
potential on that component by $R=P_g-P_g(r)$. Thus $|R|\le2$.
The test $g$ is chosen to evaluate a vector $y\in F_D$; it is not
assumed to be supported outside $D$.

We first reuse the inward increment families $u,v$ from
Lemma~\ref{thm:constrained-paired}. That proof establishes
$u/2,v/2\in K$, so their evaluations are continuous on $X_K$, and
extends $(u-v)(y)=g(y)$ from finite tails to every $y\in F_D$.

\subsection{One of two inward parts}
\label{sec:two-part-selection}

Let $u,v$ be the inward increments constructed in that proof. In addition
to the tests used there, for $0\le\lambda\le1/2$ we have
\begin{equation}\label{eq:normalized-inward}
 \frac{f+\lambda u}{\sqrt{1+\lambda^2}},\qquad
 \frac{f+\lambda v}{\sqrt{1+\lambda^2}}\ \in K.
\end{equation}
Indeed $c+\lambda U$ and $c+\lambda V$ already lie in $[-1,1]$.
Each quadratic sum is at most $1+\lambda^2$, because the supports of
the head and tail increments are disjoint. Dividing enforces the
quadratic bound and preserves the path bound and support requirements.

The identity $(u-v)(y)=g(y)$ shows that one of these two inward parts
retains at least half the evaluation:
\[
 \max\{|u(y)|,|v(y)|\}\ge |g(y)|/2\qquad(y\in F_D).
\]
This is the single-part choice, available in all four models. To see two
quantitative specializations,
suppose $\|x\|_K=1$, $x$ is supported in $D$, $f(x)=1$,
$y\in F_D$, and $\|y\|_K\ge1$. Evaluation of
\eqref{eq:normalized-inward} at the favorable endpoint gives
\[
 \max_{\pm}\|x\pm ty\|_K
 \ge\frac{1+\lambda t|g(y)|/2}{\sqrt{1+\lambda^2}}.
\]
For $0<t<1$, set $\lambda=t/2$ and take the supremum over $g$.
The maximum is at least $\sqrt{1+t^2/4}$; each endpoint is at least
one by the head test. Thus this selection yields average gain
$D(t)=(\sqrt{1+t^2/4}-1)/2$ at finite centers. Approximation of a
unit center by a finite unit center within $D(t)/2$ gives the gain
$D(t)/2$. This proves the rootless sibling estimate directly using a
single chosen tail part.

The square-root normalization can also be replaced by the test
$\sqrt{1-\theta^2}f+\theta q$, where $q$ is a chosen member of
$\{u,v\}$. If $|g(y)|\ge\|y\|_K/2$, then
$|q(y)|\ge\|y\|_K/4$. Taking $\theta=t/8$ gives
\[
 \max_{\pm}\|x\pm ty\|_K
 \ge\sqrt{1-\theta^2}+t\theta\|y\|_K/4
 \ge1+t^2/64.
\]
Both endpoints remain at least one. In the globally constrained space
this gives finite-center average gain $t^2/128$, and approximation
within $t^2/256$ gives the completed-space gain $t^2/256$.
The continuity established above applies also when the global test has
infinite support.

\subsection{Four pieces classified by attachment sign}
\label{sec:four-piece-selection}

For each $i\in\{-1,1\}$, consider those exit components with
$i=1$ when $c\ge0$ and $i=-1$ otherwise. On these components take
the potentials $R_+$ and $(-R)_+$, and take zero on all other
components and on $D$. This produces four nonnegative potentials
$G_{i,+},G_{i,-}$, all bounded by two. Their increments
$g_{i,+},g_{i,-}$ satisfy coordinatewise
\[
 |g_{i,\pm}(s)|\le |g(s)|,
 \qquad g\mathbf1_{T\setminus D}
       =\sum_{i=\pm1}(g_{i,+}-g_{i,-}).
\]
The coefficient bounds hold also on exit edges because the relative
potential there starts at zero. Consequently $g_{i,\pm}/2\in K$.
The displayed coefficient identity holds on finite tails and extends
to every $y\in F_D$ by continuity. One of the four pieces therefore
has evaluation of magnitude at least $|g(y)|/4$. In particular, when
$|g(y)|\ge\|y\|_K/2$, it retains at least $\|y\|_K/8$.

This four-piece choice is also available in all four models. In the
root-inclusive sibling model, it gives a direct argument at
an arbitrary unit center $x$. Fix $0<t<1$, put
\[
 \lambda=t/16,\qquad a=\sqrt{1-\lambda^2},\qquad
 \eta=t^2/2048,
\]
and choose a finite $x_0$ with $\|x-x_0\|\le\eta$.
Take $D$ containing its support and a test $f$ supported in $D$ with
$f(x_0)\ge\|x_0\|-\eta\ge1-2\eta$.
Choose a norming functional $\phi$ for $x$, and impose in addition
$\phi(y)=0$ on the tail space. This is still a closed
finite-codimensional subspace. For $y$ in this subspace with $\|y\|\ge1$,
choose $g\in K_{\rm p}$ with $|g(y)|\ge\|y\|/2$. Select a four-piece increment $q$ as
above and multiply it by the sign that points inward from its
attachment class. Then $a f+\lambda q$ is admissible: its head and
tail quadratic budgets are $a^2$ and $\lambda^2$, and its tail
potential changes the inherited value inward by at most $2\lambda$.
This test evaluates $x$ at least at $a-3\eta$, since it has norm at
most one and agrees with $af$ on $x_0$. Hence
\[
 \max_{\pm}\|x\pm ty\|
 \ge 1-\lambda^2-3\eta+t\lambda/8
 =1+5t^2/2048\ge1+t^2/512.
\]
The functional $\phi$ gives the lower bound one for both endpoints.
Their average therefore has gain at least $t^2/1024$.

There is a different useful normalization of the same four-piece
selection in the global model. For a unit $x\in X_{K_{\rm g}}$, put
$\alpha=t/64$ and $\eta=\alpha^2/10$; approximate $x$ by a vector
$x_0$ supported in $D$ within $\eta/10$. A head norming test for
$x_0$ has $f(x)\ge1-\eta$. For $y\in F_D$ with $\|y\|\ge1$, choose an inward piece $q$ with
$|q(y)|\ge1/8$. Its functional norm is at most two, so
$|q(x)|\le2\eta/10$. The test
\[
 H=\frac{f+\alpha q}{\sqrt{1+\alpha^2}}
\]
is admissible by disjoint global square sums and the inward path
bound. Both endpoint norms are at least $1-\eta$, and
\begin{equation}\label{eq:global-four-piece-estimate}
 \max_{\pm}\|x\pm ty\|
 \ge\frac{1-\eta-2\alpha\eta/10+\alpha t/8}
              {\sqrt{1+\alpha^2}}
 \ge1+2\alpha^2.
\end{equation}
For the last inequality, $\alpha t/8=8\alpha^2$ and
$\sqrt{1+\alpha^2}\le1+\alpha^2/2$. The numerator is
$1+(79/10)\alpha^2-\alpha^3/50$, which dominates
$(1+2\alpha^2)(1+\alpha^2/2)$ for $0\le\alpha\le1/64$.
The two endpoint bounds yield average gain at least
$19\alpha^2/20=19t^2/81920$.

\subsection{Independent component selection}
\label{sec:component-selection}

For $K=K_{\rm a}$ every test has finite support. Thus only finitely many
exit components carry nonzero increments, and we can choose pieces
component by component. We first describe the two choices and their
evaluation guarantees, then give their numerical specializations.

For the first choice, take $g\in K_{\rm a}$ and split its relative
potential as
\[
 R=2(R_+/2-(-R)_+/2).
\]
The two increment families $g^+,g^-$ have coefficients bounded by
$|g_s|/2$ and component potentials in $[0,1]$, so both belong to
$K_{\rm a}$. All component sums below are finite. The coefficient identity
$g\mathbf1_{T\setminus D}=2(g^+-g^-)$ shows that one family, denoted
$k$, satisfies $|k(y)|\ge |g(y)|/4$. Let $q_j$ be its evaluation on
the $j$th exit component. Then $\sum_j|q_j|\ge |g(y)|/4$.
Choose on each component a sign $d_j$ pointing inward from its inherited
head value. For $0<t<1$, put $r=t/32$ and $s=\sqrt{1-r^2}$. Adjoining either
$rd_jk$ or zero on each component to $sf$ gives a test. The path value
moves inward by at most $r$, and every antichain receives square sum at
most $s^2+r^2=1$. Retain the components with $d_jq_j>0$ for the plus
endpoint and those with $d_jq_j<0$ for the minus endpoint. For a vector
$x_0$ supported in $D$, the resulting pair of tests gives
\[
 \frac{\|x_0+ty\|+\|x_0-ty\|}{2}
 \ge s f(x_0)+\frac{tr}{2}\sum_j|q_j|
 \ge s f(x_0)+tr|g(y)|/8.
\]

For the second choice, split the relative potential of $g$ into its
positive and negative parts on each exit component. Select the part whose
evaluation has magnitude at least half the magnitude of the full
component evaluation, and orient it inward. If the selected
increments are $h_j$, then
\[
 \sum_j|h_j(y)|\ge |g(y)|/2.
\]
Retain either all positively evaluating pieces or all negatively
evaluating pieces, according to which group has greater total
magnitude. Their sum $h$ satisfies $|h(y)|\ge|g(y)|/4$,
$|h_s|\le|g_s|$, and $h/2\in K_{\rm a}$. All these arrays are finite.
For $0\le\alpha\le1/2$, the single test
$\sqrt{1-\alpha^2}f+\alpha h$ is admissible by the same antichain
budget and inward path argument.

We now recover the two completed-space estimates. First fix $0<t<1$
and a unit $x$ in $X_{K_{\rm a}}$. Set $r=t/32$,
$s=\sqrt{1-r^2}$, and $\eta=t^2/10000$. Choose a finite $x_0$ within
$\eta$ of $x$ and a test $f$ supported in a finite initial $D$ with
$f(x_0)\ge1-2\eta$. For $y\in F_D$ of norm at least one choose
$g\in K_{\rm a}$ with $g(y)\ge1/2$. The first choice gives
\[
 \frac{\|x_0+ty\|+\|x_0-ty\|}{2}
 \ge s(1-2\eta)+tr/16.
\]
Returning to $x$ loses at most $\eta$. Since $s\ge1-r^2$, the final
excess is at least
\begin{equation}\label{eq:antichain-component-bound}
 tr/16-r^2-3\eta=\frac{433}{640000}t^2>0.
\end{equation}
The subspace was fixed before $y$, as required by the modulus.

For the single-test specialization put $\alpha=t/16$ and
$\varepsilon=t^2/2048$. Choose finite $x_0$ within $\varepsilon$
of $x$ and a finite test $f$ with $f(x)>1-\varepsilon$; enlarge $D$
to contain both supports. For $y\in F_D$ of norm at least one choose
$g\in K_{\rm a}$ with $g(y)>\|y\|/2$. The second choice gives an
inward $h$ with $|h(y)|\ge1/8$. Since $h$ vanishes on $D$ and has
functional norm at most two, $|h(x)|\le2\varepsilon$. Evaluating
$\sqrt{1-\alpha^2}f+\alpha h$ at a favorable endpoint bounds it below by
\[
 1-\alpha^2-(1+2\alpha)\varepsilon+\alpha t/8,
\]
while the head test bounds the other by $1-\varepsilon$.
The average excess is consequently at least
\[
 \frac{\alpha t/8-\alpha^2-(2+2\alpha)\varepsilon}{2}
 \ge\frac{5t^2}{4096}\ge\frac{t^2}{1024}.
\]
This last choice keeps the full coefficient domination of $g$ and
uses only one inward test at the favorable endpoint. The common
paired proof and these selection proofs therefore provide different
witnesses under the same explicitly stated antichain constraint.
